\section{Conclusão}
\label{sec:conclusion}

Replicamos o \modus\ a partir do texto do artigo e implementamos as melhorias propostas na primeira
etapa do trabalho. O percurso produziu três resultados.

O primeiro é um defeito de avaliação. O \texttt{HSTrees} do MOA usa o atributo de classe como
dimensão de corte, e o uso padrão do CapyMOA entrega o rótulo verdadeiro ao detector. No
ASEADOS-SDN-IoT, o F1 de 99,73\% da nossa réplica cai para 0,00\% quando o rótulo sai do fluxo. O
efeito depende da base, o que o torna fácil de não ver.

O segundo é um defeito de projeto. Sem o vazamento, o treino seletivo do \modus\ e da sua versão
normalizada deixa de 99,96\% a 100\% do ataque entrar no modelo nas três bases. O \emph{flood}
vira o perfil normal, e a AUC do score fica abaixo de 50\%. Onde o \modus\ original ainda detecta,
é pela raridade do protocolo do ataque, e não pelo volume.

O terceiro é o \modusp, com quatro defesas causais e sem rótulos. Com uma única configuração, o F1
vai de 0,08\% para 62,99\% no ASEADOS e de 76,66\% para 91,17\% no Edge-IIoTset, onde o método
passa a detectar os \emph{floods} TCP e HTTP que o original perdia. No CTU-13 as defesas reduzem o
envenenamento de 100\% para 2,41\%, mas o F1 a priori é 0,01\%, porque o score do ataque fica
milésimos abaixo do limiar. Ataques não-DDoS continuam, em sua maioria, sem detecção.

As limitações apontam os próximos passos. A calibração do limiar sem rótulos é a principal: nas
duas bases em que medimos, o modelo congelado separa as classes com AUC de 94\% e a regra de quantil
desperdiça a separação. A decisão por janela e por destino é o caminho que consideramos mais
promissor. Faltam ainda a renovação periódica do modelo congelado, a avaliação com mais sementes
nas bases grandes e bases com ataques curtos ou de baixa taxa.

Para quem avalia detectores de anomalia em fluxos, ficam três recomendações práticas: verificar
que o resultado não muda com rótulos constantes; reportar a fração do ataque usada no treino, e
não só o F1; e fixar uma configuração única para todas as bases antes de olhar os resultados. O
código, os \emph{logs} e as métricas de todas as rodadas estão no repositório do trabalho.
